\section{Proof of \autoref{alg:method} Optimality and Approximate Optimality}
\label{sec:optimality-proofs}

In \autoref{sec:relaxation}, we claim that \autoref{alg:method} produces
optimal non-discriminatory classifiers in exact calibration scenarios,
and near-optimal classifiers in approximate calibration scenarios.

We begin with classifiers $h_1$ and $h_2$ be classifiers for groups $G_1$ and
$G_2$, with calibrations $\epsilon(h_1) \le \delta_{cal}$ and $\epsilon(h_2)
\le \delta_{cal}$. As before, assume that we cannot strictly
improve the cost of either $h_1$ or $h_2$ without worsening calibration: i.e. $h_1$ and
$h_2$ is $\fairconstname$ and calibration. We will now show that
\autoref{alg:method} produces classifiers that are
near-optimal with respect to both the false-positive and false-negative rates
among calibrated classifiers satisfying the equal-cost constraint.

First, we show that interpolation preserves approximate calibration:
\begin{thm}[Approximate Optimality of \autoref{alg:method}]
  Given $\fair h_2$, which is the classifier produced by \autoref{alg:method},
  we have that $\calib{\fair h_2} \leq \left( 1 - \alpha \right) \calib{h_2}$, where
  $\alpha \in [0, 1]$ is the interpolation parameter in \eqref{eqn:interpolation}.
  \label{thm:calibration}
\end{thm}
%
\begin{proof}
  %Again, let $h_t$ be a classifier for group $G_t$, and let $h^{\br_t}$ be the
  %trivial classifier for $G_t$.
  %Finally, we assume that $\fair h_t$ randomizes between the two classifiers, i.e.
  %\[
    %\fair h_t(\x) = (1-B) h_t(\x) + B h^{\br_t}(\x)
  %\]
  %where $B$ is a Bernoulli random variable with parameter $\alpha$.
  We can calculate the calibration of $\fair h_2(\x)$ as follows:
  %
  \begin{align*}
    \calib{\fair h_2}
    &= \mathop{\ev}_{B, G_2} \Bigl\lvert \Pr_{G_2} \bigl[ y\!=\!1 \mid \fair
    h_2(\x)\!=\!p \bigr] - p \Bigr\rvert
    \\
    &= \int_{0}^1 \Bigl\lvert \Pr_{G_2} \bigl[ y\!=\!1 \mid \fair
    h_2(\x) \!=\!p \bigr] - p \Bigr\rvert
    %\\ &\phantom{=} \times
    \Pr_{G_2} \bigl[ \fair h_2(\x)\!=\!p\bigr] dp
    %\\
    %&= \int_{0}^1 \Biggl( \Bigl\lvert \Pr_{G_2} \bigl[ y\!=\!1 \mid h_2(\x)
    %\!=\!p \bigr] - p \Bigr\rvert \Pr_{G_2} \bigl[ h_2(\x)\!=\!p, B\!=\!1 \bigr]
    %%\\ &\phantom{=}
    %+ \Bigl\lvert \Pr_{G_2} \bigl[ y\!=\!1 \mid h^{\br_2}(\x) \!=\!p
    %\bigr] - p \Bigr\rvert \Pr_{G_2} \bigl[ h^{\br_2}(\x)\!=\!p, B\!=\!0 \bigr] \Biggr) dp
    %\\
    %&= \int_{0}^1 \Biggl( \alpha \Bigl\lvert \Pr_{G_2} \bigl[ y\!=\!1 \mid
    %h_2(\x) \!=\!p \bigr] - p \Bigr\rvert \Pr_{G_2} \bigl[ h_2(\x)\!=\!p \bigr]
    %%\\ &\phantom{=}
    %+  \bigl( 1 \! - \! \alpha \bigr) \Bigl\lvert \Pr_{G_2}
    %\bigl[ y\!=\!1 \mid h^{\br_2}(\x) \!=\!p \bigr] - p \Bigr\rvert \Pr_{G_2}
    %\bigl[ h^{\br_2}(\x)\!=\!p \bigr] \Biggr) dp
    %\\
    %&= (1-\alpha) \calib{h_2} + \alpha \calib{h^{\br_2}}.
  \end{align*}
  For $p \ne \mu_2$, $|\Pr_{G_2} \bigl[ y\!=\!1 \mid \fair h_2(\x) \!=\!p \bigr]
  - p| = |\Pr_{G_2} \bigl[ y\!=\!1 \mid h_2(\x) \!=\!p \bigr] - p|$ and
  $\Pr_{G_2} \bigl[ \fair h_2(\x)\!=\!p\bigr] = (1-\alpha)\Pr_{G_2} \bigl[
  h_2(\x)\!=\!p\bigr]$.

  For $p = \mu_2$, let $\beta = \Pr_{B,G_2}\bigl[B \!=\! 1 \mid \fair h_2(\x) =
  p\bigr]/\Pr_{G_2} \bigl[\fair h_2(\x) = p\bigr]$. Note that $\beta \ge
  \alpha$. Then,
  \begin{align*}
    \Pr_{G_2} \bigl[ y\!=\!1 \mid \fair h_2(\x) \!=\!p \bigr] &=
    (1-\beta) \Pr_{G_2} \bigl[ y\!=\!1 \mid h_2(\x) \!=\!p \bigr] + \beta
    \Pr_{G_2} \bigl[ y\!=\!1 \mid h^{\br_2}(\x) \!=\!p \bigr] \\
    &= (1-\beta) \Pr_{G_2} \bigl[ y\!=\!1 \mid h_2(\x) \!=\!p \bigr] + \beta p
    %&= (1-\alpha) |\Pr_{G_2} \bigl[ y\!=\!1 \mid h_2(\x) \!=\!p \bigr] - p|
  \end{align*}
  because $h^{\br_2}$ is perfectly calibrated. Moreover, note that $\Pr_{G_2}
  \bigl[ \fair h_2(\x)\!=\!p\bigr] = \Pr_{G_2} \bigl[ h_2(\x)\!=\!p\bigr] /
  (1-\beta)$
  Using this, we have
  $|\Pr_{G_2} \bigl[ y\!=\!1 \mid \fair h_2(\x) \!=\!p
  \bigr] - p | \Pr_{G_2} \bigl[ \fair h_2(\x)\!=\!p\bigr] =
  |\Pr_{G_2} \bigl[ y\!=\!1 \mid h_2(\x) \!=\!p
  \bigr] - p | \Pr_{G_2} \bigl[ h_2(\x)\!=\!p\bigr]$.
  Thus,
  \begin{align*}
    \calib{\fair h_2}
    &= \int_{0}^1 \Bigl\lvert \Pr_{G_2} \bigl[ y\!=\!1 \mid \fair
    h_2(\x) \!=\!p \bigr] - p \Bigr\rvert
    \Pr_{G_2} \bigl[ \fair h_2(\x)\!=\!p\bigr] dp \\
    &\le \int_{0}^1 \Bigl\lvert \Pr_{G_2} \bigl[ y\!=\!1 \mid
    h_2(\x) \!=\!p \bigr] - p \Bigr\rvert
    \Pr_{G_2} \bigl[ h_2(\x)\!=\!p\bigr] dp \\
    &= \calib{h_2}
  \end{align*}
\end{proof}

Next, we observe that by \autoref{lma:cost_error}, for any
classifiers $h_1'$ and $h_2'$ with $\epsilon(h_1') \le \delta_{cal}$ and
$\epsilon(h_2') \le \delta_{cal}$ satisfying the equal-cost constraint, it must be
the case that $\fp{h_t'} \ge \fp{\fair h_t} - \frac{4 \delta_{cal}}{1-\mu_t}$
and $\fn{h_t'} \ge \fp{\fair h_t} - \frac{4 \delta_{cal}}{\mu_t}$ for $t = 1, 2$.

%\begin{thm}[Approximate false-positive optimality for \autoref{alg:method}]
  %\label{thm:approx-fp-optimality}
  %Let $\fair h_1^{*fp}$ and $\fair h_2^{*fp}$ be the false-positive-optimal non-discriminatory
  %classifiers under the constraint $\epsilon(\fair h_1^{*fp}) \leq \delta_{cal}^{(1)}$
  %and $\epsilon(\fair h_2^{*fp}) \leq \delta_{cal}^{(2)}$.
  %Then we have that
  %$$\fp{h_1} - \fp{\fair h_1^{*fp}} \leq
  %\frac{4 \delta_{cal}^{(1)}}{1 - \mu_1}.$$
  %and
  %$$\fp{\fair h_2} - \fp{\fair h_2^{*fp}} \leq
    %4 b_2 \delta_{cal}^{(2)}
    %+ \left( \frac{4 a_1 \delta_{cal}^{(1)}}{(1 - \mu_1) \left( \mu_2 a_2 + (1-\mu_2)b_2 \right)} \right),$$
  %%
  %where $\fair h_2$ is the derived classifier from \autoref{alg:method}.
%\end{thm}
%%
%\begin{proof}
  %Because $h_1$ is Pareto-dominant, we have that if $\fp{\fair h_1^{*fp}}$ is less than $\fp{h_1}$,
  %then $\fn{h_2} \leq \fn{\fair h_1^{*fp}}$.
  %We can use this fact together with \autoref{lem:fpfn} to lower-bound $\fp{\fair h_1^{*fp}}$:
  %%
  %$$\fp{\fair h_1^{*fp}} - \fp{h_1} \leq \frac{4 \delta_{cal}^{(1)}}{1 - \mu_1}.$$
  %%
  %This in turn leads to a lower bound on $\fairconst{\fair h_1^{*fp}}{1}$:
  %%
  %\begin{align*}
    %\fairconst{\fair h_1^{*fp}}{1} - \fairconst{h_1}{1}
    %&= a_1 \left( \fp{\fair h_1^{*fp}} - \fp{h_1} \right)
    %+ b_1 \left( \fn{\fair h_1^{*fp}} - \fn{h_1} \right)
      %\\
      %&\leq a_1 \left( \frac{4 \delta_{cal}^{(1)}}{1 - \mu_1} \right)
      %+ b_1 \left( \fn{h_1^{*fp}} - \fn{h_1} \right)
      %\\
      %&\leq a_1 \left( \frac{4 \delta_{cal}^{(1)}}{1 - \mu_1} \right).
  %\end{align*}

  %Now we will derive a relationship between $\fairconst{\fair h_2}{2}$ and $\fp{\fair h_2}$.
  %Using \autoref{lem:approx-linear-cal}, we can generate lower and upper bounds only
  %in terms of the generalized false-positive rate:
  %%
  %\begin{align*}
    %\fairconst{\fair h_2}{2} &= a_2 \fp{\fair h_2} + b_2 \fn{\fair h_2}.
    %\\
    %&\leq \left( \frac{\mu_2 a_2 + (1-\mu_2)b_2}{\mu_2} \right) \fp{\fair h_2}
    %+ \frac{2 b_2 \delta_{cal}^{(2)}}{\mu_2}
  %\end{align*}
  %%
  %and
  %\begin{align*}
    %\fairconst{\fair h_2}{2}
    %&\geq \left( \frac{\mu_2 a_2 + (1-\mu_2)b_2}{\mu_2} \right) \fp{\fair h_2}
    %- \frac{2 b_2 \delta_{cal}^{(2)}}{\mu_2}
  %\end{align*}
  %%
  %Note that these same equations hold for $\fair h_2^{*fp}$. Using these
  %equations, we can derive an upper bound for $\fp{\fair h_2}$ and a lower
  %bound for $\fp{\fair h_2^{*fp}}$:
  %%
  %\begin{align*}
    %\fp{\fair h_2}
    %&\leq \left( \frac{\mu_2 \fairconst{\fair h_2}{2}}{\mu_2 a_2 + (1-\mu_2)b_2} \right)
    %+ 2 b_2 \delta_{cal}^{(2)}
    %\\
    %\fp{\fair h_2^{*fp}}
    %&\geq \left( \frac{\mu_2 \fairconst{\fair h_2^{*fp}}{2}}{\mu_2 a_2 + (1-\mu_2)b_2} \right)
    %- 2 b_2 \delta_{cal}^{(2)}.
  %\end{align*}
  %%
  %By the no-disparity constraint, we have that
  %$\fairconst{\fair h_2}{2} = \fairconst{h_1}{1}$. Similarly,
  %we know that $\fairconst{\fair h_2^{*fp}}{2} \geq \fairconst{\fair h_1^{*fp}}{1} \geq
  %\fairconst{h_1}{1} - a_1 \left( \frac{4 \delta_{cal}^{(1)}}{1 - \mu_1} \right)$.
  %Putting this all together, we have:
  %%
  %\begin{align*}
    %\fp{\fair h_2}
    %&\leq \left( \frac{\mu_2 \fairconst{h_1}{1}}{\mu_2 a_2 + (1-\mu_2)b_2} \right)
    %+ 2 b_2 \delta_{cal}^{(2)}
    %\\
    %&= \left( \frac{\mu_2 \fairconst{h_1}{1} - 4 a_1 \delta_{cal}^{(1)} / (1 - \mu_1)}
    %{\mu_2 a_2 + (1-\mu_2)b_2} \right)
    %+ 2 b_2 \delta_{cal}^{(2)}
    %+ \left( \frac{4 a_1 \delta_{cal}^{(1)}}{(1 - \mu_1) \left( \mu_2 a_2 + (1-\mu_2)b_2 \right)} \right)
    %\\
    %&\leq \left( \frac{\mu_2 \fairconst{\fair h_2^{*fp}}{2}}
    %{\mu_2 a_2 + (1-\mu_2)b_2} \right)
    %+ 2 b_2 \delta_{cal}^{(2)}
    %+ \left( \frac{4 a_1 \delta_{cal}^{(1)}}{(1 - \mu_1) \left( \mu_2 a_2 + (1-\mu_2)b_2 \right)} \right)
    %\\
    %&\leq \fp{\fair h_2^{*fp}}
    %+ 4 b_2 \delta_{cal}^{(2)}
    %+ \left( \frac{4 a_1 \delta_{cal}^{(1)}}{(1 - \mu_1) \left( \mu_2 a_2 + (1-\mu_2)b_2 \right)} \right)
  %\end{align*}
%\end{proof}
%%
%Similarly, we can derive a result for false-negative optimality.
%%
%\begin{thm}[Approximate false-negative optimality for \autoref{alg:method}]
  %\label{thm:approx-fn-optimality}
  %Let $\fair h_1^{*fn}$ and $\fair h_2^{*fn}$ be the false-negative-optimal non-discriminatory
  %classifiers under the constraint $\epsilon(\fair h_1^{*fn}) \leq \delta_{cal}^{(1)}$
  %and $\epsilon(\fair h_2^{*fn}) \leq \delta_{cal}^{(2)}$.
  %Then we have that
  %$$\fn{h_1} - \fn{\fair h_1^{*fp}} \leq
  %\frac{4 \delta_{cal}^{(1)}}{\mu_1}.$$
  %and
  %$$\fn{\fair h_2} - \fn{\fair h_2^{*fp}} \leq
    %4 b_2 \delta_{cal}^{(2)}
    %+ \left( \frac{4 a_1 \delta_{cal}^{(1)}}{\mu_1 \left( (1 - \mu_2) a_2 + \mu_2 b_2 \right)} \right),$$
  %%
  %where $\fair h_2$ is the derived classifier from \autoref{alg:method}.
%\end{thm}
%%
%It is important to note that both bounds degrades linearly with respect to calibration.
Thus, approximately calibrated classifiers will be approximately optimal. From
this result, it is easy to derive the optimality result for perfectly-calibrated
classifiers.
%
\begin{thm}[Exact Optimality of \autoref{alg:method}]
  \label{thm:optimality}
  \autoref{alg:method} produces the classifiers $h_1$ and $\fair h_2$ that
  satisfy both perfect calibration and the equal-cost constraint with the lowest
  possible generalized false positive and false negative rates.
  %Let $h_1$ and $h_2$ be (discriminatory) classifiers which achieve the
  %minimum false-positive and false-negative rates for groups $G_1$ and $G_2$, respectively
  %(i.e. cannot be strictly improved).
  %If $h_1$ and $h_2$ are perfectly calibrated, then \autoref{alg:method} produces
  %the optimal perfectly calibrated non-discriminatory classifiers with respect
  %to false-positive/false-negative rates.
\end{thm}
